\label{sec:alabama}

\subsection{The Hamilton Method}

\begin{definition}[Hamilton Method]
  Given populations $p_1, \ldots, p_n$ and House size $H$:
  \begin{enumerate}
    \item Compute exact quotas $q_i = p_i \times H / N$ where $N = \sum_j p_j$.
    \item Let $r = H - \sum_i \lfloor q_i \rfloor$ be the number of
      ``remainder seats'' (necessarily $0 \leq r < n$).
    \item Assign lower quota $\lfloor q_i \rfloor$ to each state.
    \item Distribute the $r$ remainder seats to the $r$ states with
      largest fractional parts $f_i = q_i - \lfloor q_i \rfloor$,
      breaking ties by population.
  \end{enumerate}
\end{definition}

Hamilton satisfies the quota rule by construction:
every state receives either $\lfloor q_i \rfloor$ or $\lceil q_i \rceil$
seats.
This is its principal advantage over divisor methods.

\subsection{Definition of the Alabama Paradox}

\begin{definition}[Alabama Paradox]
  An apportionment method $M$ satisfies the Alabama Paradox if there
  exist populations $p_1, \ldots, p_n$ and House size $H$ such that
  $M(\vec{p}, H)$ gives state $i$ seats $s_i$, but
  $M(\vec{p}, H+1)$ gives state $i$ seats $s_i - 1 < s_i$.
  (Increasing the House size causes a state to lose a seat.)
\end{definition}

\subsection{The 1880 Alabama Paradox: Historical Example}

The 1880 Census populations for the relevant states were:
\begin{center}
\begin{tabular}{lr}
\toprule
State & 1880 Population \\
\midrule
Alabama       & 1,262,505 \\
Texas         & 1,591,749 \\
Illinois      & 3,077,871 \\
Total (all states) & 50,155,783 \\
\bottomrule
\end{tabular}
\end{center}

At $H = 299$, Alabama's exact quota is:
$q_{\text{AL}} = 1{,}262{,}505 \times 299 / 50{,}155{,}783 \approx 7.530$.
Lower quota: 7. Fractional part: 0.530.
Alabama receives 8 seats (lower quota 7 + 1 remainder seat,
as its 0.530 fractional part ranks among the top $r$ remainder receivers).

At $H = 300$, Alabama's exact quota is:
$q_{\text{AL}} = 1{,}262{,}505 \times 300 / 50{,}155{,}783 \approx 7.545$.
Lower quota: 7. Fractional part: 0.545.
One might expect Alabama to retain its remainder seat (its fractional part increased).
However, Texas's exact quota at $H = 300$ is:
$q_{\text{TX}} = 1{,}591{,}749 \times 300 / 50{,}155{,}783 \approx 9.516$.
Lower quota: 9. Fractional part: 0.516.
At $H = 299$, Texas had fractional part $0.516 \times 299/300 \approx 0.515$.
Illinois at $H = 300$ has fractional part approximately 0.549.

The critical change at $H = 300$ is that the addition of one seat
increases the lower quotas for multiple states, reducing the number
of remainder seats $r$, and changes the fractional-part ranking.
Alabama's fractional part of 0.545 is now \emph{below} the fractional
parts of Illinois (0.549) and possibly other states, so Alabama
loses its remainder seat.
Result: Alabama receives 7 seats at $H = 300$.

\begin{example}[Synthetic Alabama Paradox]
  Consider 3 states with populations $(6, 6, 2)$ and $H$ varying:
  \[
    N = 14, \quad q_i = \frac{p_i \cdot H}{14}.
  \]
  At $H = 7$: quotas $(3, 3, 1)$ exactly, no fractional parts,
  Hamilton gives $(3, 3, 1)$.
  At $H = 8$: quotas $(3.43, 3.43, 1.14)$, lower quotas $(3, 3, 1)$,
  $r = 1$ remainder, fractional parts $(0.43, 0.43, 0.14)$,
  tie broken by population: state 1 gets remainder.
  Hamilton gives $(4, 3, 1)$.
  At $H = 14$: quotas $(6, 6, 2)$, Hamilton gives $(6, 6, 2)$.

  Now consider $H = 9$: quotas $(3.86, 3.86, 1.29)$, lower quotas $(3, 3, 1)$,
  $r = 2$ remainders.
  Hamilton gives $(4, 4, 1)$.
  State 3 had 1 seat at $H=7$, 1 at $H=8$, 1 at $H=9$ ---
  no paradox here. But with different population ratios, a paradox occurs.

  For a clean 3-state example exhibiting the paradox:
  populations $(5, 5, 2)$, $N = 12$:
  \begin{itemize}
    \item $H = 11$: quotas $(4.583, 4.583, 1.833)$, lower quotas $(4,4,1)$,
      $r=2$, fractional parts $(0.583, 0.583, 0.833)$.
      State 3 gets one remainder, state 1 gets the other.
      Hamilton: $(5, 4, 2)$.
    \item $H = 12$: quotas $(5, 5, 2)$, lower quotas $(5,5,2)$, $r=0$.
      Hamilton: $(5, 5, 2)$.
    \item $H = 13$: quotas $(5.417, 5.417, 2.167)$, lower quotas $(5,5,2)$,
      $r=1$, fractional parts $(0.417, 0.417, 0.167)$.
      State 1 gets the remainder (tie with state 2 broken by label).
      Hamilton: $(6, 5, 2)$.
  \end{itemize}
  State 2 had 4 seats at $H=11$ but 5 at $H=12$: no paradox.
  Constructing a clean Alabama paradox requires populations that create
  the fractional-part crossing \citep{balinski1982fair}.
\end{example}

\subsection{Why Hamilton Is Susceptible}

The Alabama Paradox occurs because the number of remainder seats
$r = H - \sum_i \lfloor q_i \rfloor$ does not increase monotonically
with $H$.
When $H$ increases by 1:
\begin{itemize}
  \item Every exact quota $q_i = p_i \times H / N$ increases by $p_i/N$.
  \item Some lower quotas $\lfloor q_i \rfloor$ increase by 1
    (when $q_i$ crosses an integer).
  \item The total increase in lower quotas equals
    $|\{i : q_i \text{ crosses an integer at this } H\}|$,
    which can be 0, 1, 2, or more.
\end{itemize}
If the total lower quotas increase by 2 when $H$ increases by 1,
then $r = H - \sum \lfloor q_i \rfloor$ \emph{decreases} by 1.
This means one fewer remainder seat is available, and some state
that had a remainder seat loses it.
The fractional-part ranking then determines which state loses the seat.
If the losing state is not the one whose quota crossed an integer
(i.e., not the state that ``earned'' the lower-quota increase),
the result is counterintuitive: a state unrelated to the quota-crossing
event loses a seat.
